% IDENTITY-REWRITE: new Limitations section

\section{Limitations}
\label{sec:limitations}

Several limitations bound what this paper can claim.

\paragraph{Within-market selection.}
The main threat to interpretation. Item, year, and PBU fixed effects
absorb time-invariant cross-market heterogeneity, and the cross-fit
removes the mechanical link between the FL classification and prices.
Neither resolves the possibility that unobserved, time-varying market
characteristics attract both FL firms and higher prices within the
same item--year--PBU cell. The sensitivity analysis ($RV = 17.5\%$)
bounds how large such a confounder would need to be, but bounding is
not identification.

\paragraph{Pre-trends in the event study.}
Positive coefficients at $t = -2$ and $t = -3$ before first FL entry
are equally consistent with strategic selection by cartels into
high-value markets and with within-market confounding that predates
FL arrival. This pattern precludes a causal interpretation of the
event study and is the primary reason the paper is positioned as
diagnostic rather than causal.

\paragraph{FL persistence.}
Only 10\% of firms classified as FL in 2009--2013 reappear as FL in
2014--2019. This limits the interpretation of FL as a stable firm
type. It is more accurately read as an environment-level marker
subject to cartel reorganization, firm exit, or classification noise.

\paragraph{Zero-win sensitivity.}
Relaxing the zero-win condition to include firms with 1--2\% win
rates weakens the price coefficient substantially. The screen depends
on a sharp boundary that sophisticated cartels could evade by allowing
cover bidders an occasional win.

\paragraph{CADE overlap.}
Conditional on participation intensity, FL firms are indistinguishable
from other always-losers in their overlap with CADE-convicted
cartelists ($p = 0.93$). The screen captures market-level exposure to
cartel-affected environments, not firm-level cartel membership.

\paragraph{Structural model.}
The model calibrates cover-bid distributions and selects between
regimes, but does not identify the markup or the causal price effect.
Its welfare implications are illustrative, not causal
(Appendix~\ref{app:extensions}).

\paragraph{External validity.}
All results come from S\~ao Paulo's BEC platform. Generalization to
federal procurement, other Brazilian states, or other countries
requires replication with the same FL definition applied to different
institutional settings.
